
Current RLHF optimizes language models for helpfulness alone, neglecting the $Human\rightarrow AI$ direction of alignment: controllability. This work introduces bidirectional alignment, augmenting standard helpfulness rewards with IFEval-derived controllability signals via the combined reward R = $\alpha \cdot R_helpfulness + \beta \cdot R_IFEval$. Discrete constraint checks are converted into differentiable rewards using sigmoid soft thresholds.

Proof-of-concept experiments on Llama-3-8B-Instruct validate three predictions: (1) bidirectional models achieve +3.4 percentage points (pp) held-out IFEval strict accuracy over helpfulness-only baselines; (2) at $\alpha =0.8$, models retain 96.4\% of baseline AlpacaEval performance; (3) explicit constraint training correlates with implicit safety improvements, with +2.35pp TruthfulQA MC1 and +1.54pp BBQ improvements over baseline maxima.

The correlation between IFEval gains and safety gains (r=0.85-0.94) suggests that explicit constraint training may build general constraint-following capacity. This work provides an initial empirical test of the bidirectional alignment framework proposed by Sun et al. (2024), introduces IFEval as a differentiable training signal, and characterizes the helpfulness-controllability trade-off at proof-of-concept scale.
